\begin{textsample}{NEG Dim 1 – 2018\_c5 – Score 12.00 – 2018\_c5/enade\_2018\_ci\_cont\_3fded95eb1.txt}  \label{ex:f1_neg_371}
CAPÍTULO 5 DISTRIBUIÇÃO DOS CONCEITOS 5.1 PANORAMA NACIONAL DA DISTRIBUIÇÃO DOS CONCEITOS O Conceito Enade é um indicador de qualidade que avalia os cursos por intermédio dos desempenhos dos estudantes a partir dos resultados do Enade.
Os valores possíveis vão de 1 ( pior situação ) a 5 ( melhor situação ).
Os cursos com apenas um ou sem nenhum concluinte participante não obtêm o Conceito Enade, ficando Sem Conceito ( SC ).
Na Tabela 5.1, são apresentadas a quantidade e a distribuição de cursos de Ciências Contábeis participantes do Enade/2018, por faixa de conceito e Grande Região.
Enfatiza -se, mais uma vez, que a diferença entre os cursos tabulados neste capítulo e no capítulo 2 corresponde aos cursos Sem Conceito, ou seja àqueles, em princípio, sem alunos concluintes que houvessem participado da prova ou àqueles que tivessem tão somente um aluno concluinte¹⁹.
Observando -se os dados da Tabela 5.1, nota -se que, dos 1.101 cursos participantes, 478 ( 43,4 \% ) classificaram -se com conceito 3, o valor modal.
Esse foi, também, o conceito modal nas regiões Sudeste ( 45,2 \% ) e Sul ( 53,4 \% ).
O conceito 2 foi o segundo mais frequente em nível nacional ( 31,6 \%, correspondendo a 348 cursos ), e o conceito 4, o terceiro ( 14,5 \%, correspondendo a 160 cursos ).
Houve, ainda, 50 cursos ( 4,5 \% ) que receberam conceito 1, e 42 cursos ( 3,8 \% ) que receberam conceito 5.
Dos 1.101 cursos de Ciências Contábeis, 23 ( 2,1 \% ) ficaram Sem Conceito ( SC ).
A região Norte participou com 77 cursos ou 7,0 \% do total nacional.
Desses, 38 cursos receberam o conceito modal 2, o que equivale a 49,4 \% do total regional.
A 27 cursos ( 35,1 \% ) ¹⁹ Esses cursos com somente um concluinte não foram considerados no capítulo 2. atribuiu -se conceito 3, e a outros seis cursos, os conceitos 1 e 4 ( 7,8 \%, \textbf{cada} ).
Nessa região, nenhum curso ficou Sem Conceito ( SC ) ou recebeu conceito 5.
A região Nordeste participou com 201 cursos ou 18,3 \% do total nacional.
Nessa região, 86 cursos ( 42,8 \% em termos regionais ) obtiveram conceito 2, conceito modal para a região.
Receberam os conceitos 3 e 4, respectivamente, 74 cursos ( 36,8 \% ) e 17 cursos ( 8,5 \% ).
Ainda nessa região, seis cursos ( 3,0 \% ) foram avaliados com conceito 5, e outros 12 cursos ( 6,0 \% ) com conceito 1.
Seis cursos ficaram Sem Conceito ( SC ) na região Nordeste.
Dos 462 cursos participantes da região Sudeste ( 42,0 \% do total nacional ), o conceito 3 foi o conceito modal, como já anteriormente assinalado, tendo sido obtido por 209 cursos ( 45,2 \% ).
O conceito 1 foi atribuído a 14 cursos ( 3,0 \% ), e o conceito 2, a 119 cursos ( 25,8 \% ).
Receberam os conceitos 4 e 5, respectivamente, 81 cursos ( 17,5 \% ) e 28 cursos ( 6,1 \% ).
Nessa região, 11 cursos ficaram Sem Conceito ( SC ).
Os 236 cursos da região Sul corresponderam a 21,4 \% do total nacional.
Como já indicado, houve predominância do conceito 3, atribuído a 53,4 \% dos cursos da região ( conceito modal ), ou seja, atribuído a 126 dos 236 cursos participantes na região Sul.
O conceito 2 foi atribuído a 56 cursos ( 23,7 \% ), e os conceitos 4 e 5, respectivamente, a 42 cursos ( 17,8 \% ) e a cinco cursos ( 2,1 \% ).
Três cursos ficaram Sem Conceito ( SC ), e quatro cursos ( 1,7 \% ) receberam conceito 1.
Já dos 125 cursos participantes na região Centro-Oeste ( 11,4 \% do total nacional ), 49 cursos ( 39,2 \% em termos regionais ) receberam conceito 2, o conceito modal.
A 14 cursos ( 11,2 \% ) foram atribuídos conceito 1 e a 42 cursos ( 33,6 \% ), conceito 3.
A 14 cursos ( 11,2 \% ) o conceito 4, a três cursos ( 2,4 \% ), o conceito 5.
Nessa região, três cursos ficaram Sem Conceito ( SC ).
No Gráfico 5.1, é apresentada a distribuição cumulativa do Conceito Enade segundo Grande Região.
Em um gráfico de distribuição cumulativa, em \textbf{cada} valor de abscissa ( eixo x ), o valor apresentado no eixo das ordenadas ( eixo y ) é o valor acumulado da variável.
Por exemplo, na região Norte ( linha azul ), todos os cursos têm conceito 4 ou abaixo disso, já que a poligonal cumulativa alcança 100,0 \% para o conceito 4.
A região Centro-Oeste ( linha laranja ) e a região Nordeste ( linha verde ) apresentam uma distribuição semelhante entre si, com um pouco menos de 90,0 \% dos cursos com conceito 3 ou abaixo desse valor.
Em linhas gerais, em um gráfico de distribuição cumulativa, poligonais mais à esquerda representam Grandes Regiões com pior distribuição, e poligonais mais à direita, Grandes Regiões com melhores conceitos.
A região Norte ( linha azul ) apresenta os cursos com conceitos mais baixos em relação as demais regiões, ao passo que a região Sudeste ( linha vermelha ), seguida da região Sul ( linha rosa ), apresentam uma distribuição com valores maiores.
Já a 5.2 CONCEITOS POR CATEGORIA ADMINISTRATIVA, POR MODALIDADE DE ENSINO E GRANDE REGIÃO Na Tabela 5.2, é apresentada a distribuição dos cursos de Ciências Contábeis participantes do Enade/2018, por Categoria Administrativa e Modalidade de Ensino, de acordo com os conceitos por eles alcançados, segundo as Grandes Regiões brasileiras.
Dos 1.101 cursos participantes, 145 ( 13,2 \% ) eram ministrados em Instituições Públicas, e 956 ( 86,8 \% ), em Instituições Privadas.
Quanto à Modalidade de Ensino, 48 cursos eram a Distância, e 1.053, Presenciais.
De acordo com as informações da Tabela 5.2, em termos nacionais, entre Instituições Públicas e Privadas, observa -se que, dos 42 cursos avaliados com conceito 5, 15 eram oferecidos em IES Públicas, e 27, em IES Privadas.
Dos 145 cursos participantes de IES Públicas, o conceito 3 foi o valor modal, atribuído a 56 cursos ( 38,6 \% ).
Dos demais cursos de Instituições Públicas participantes, sete cursos ( 4,8 \% ) foram avaliados com conceito 1, 28 cursos ( 19,3 \% ), com conceito 2, 38 cursos ( 26,2 \% ), com conceito, 4 e 15 cursos, com conceito 5.
Um curso ficou Sem Conceito ( SC ) nesta categoria.
Na rede privada, o conceito modal também foi 3, atribuído a 422 cursos, o correspondente a 44,1 \% dos 956 cursos da categoria.
Dos demais cursos participantes, 43 cursos ( 4,5 \% ) receberam conceito 1, 320 cursos ( 33,5 \% ), conceito 2, 122 cursos ( 12,8 \% ), conceito 4 e 27 cursos ( 2,8 \% ), conceito 5.
Nesta Categoria Administrativa, 22 cursos ficaram Sem Conceito ( SC ).
Dos 1.053 cursos participantes da Modalidade de Ensino Presencial, o conceito 3 foi o valor modal, atribuído a 459 cursos ( 43,6 \% ).
Dos demais cursos nesta Modalidade de Ensino, 50 cursos ( 4,7 \% ) receberam conceito 1, 330 cursos ( 31,3 \% ) receberam conceito 2, 151 cursos ( 14,3 \% ), conceito 4, e 41 cursos ( 3,9 \% ), conceito 5.
Nesta Modalidade de Ensino, 22 cursos ficaram Sem Conceito ( SC ).
Dos 48 cursos a Distância informados, um curso ficou Sem Conceito ( SC ), e nenhum curso recebeu o conceito 1.
O conceito 2 foi atribuído a 18 cursos ( 37,5 \% ) e o conceito 3, a 19 cursos ( 39,6 \% ).
Nove cursos ( 18,8 \% ) receberam conceito 4 e um curso, conceito 5.
Na análise por região, observa -se que, na região Norte, as Instituições Privadas participaram com 61 cursos ( 79,2 \% do total regional ), nenhum dos quais ficou Sem Conceito ( SC ) ou recebeu conceito 5.
O conceito modal para as Instituições Privadas na região foi 2, atribuído a 33 cursos, correspondendo a 54,1 \% dessa categoria na região.
Foram 20 cursos avaliados com conceito 3 e três cursos, receberam conceito 4.
Cinco cursos receberam conceito 1.
As Instituições Públicas participaram com 16 cursos na região Norte ( 20,8 \% do total regional ), aos quais foram atribuídos os conceitos 1 ( um curso ), 2 ( cinco cursos ), 3 ( sete cursos, conceito modal ) e 4 ( três cursos ).
Nenhum curso ficou Sem Conceito ( SC ) ou recebeu conceito 5.
Dos 77 cursos oferecidos na região Norte, 76 ( 98,7 \% ) eram na Modalidade de Ensino Presencial e um ( 1,3 \% ), na Modalidade a Distância.
Dos 76 cursos oferecidos na Modalidade Presencial, nenhum curso ficou Sem Conceito ( SC ) ou recebeu conceito 5.
Os demais cursos obtiveram os conceitos 1 ( seis cursos ), 2 ( 37 cursos, conceito modal ), 3 ( 27 cursos ) e 4 ( seis cursos ).
Na Modalidade a Distância, um curso recebeu conceito 2 ( 100 \% do total da Modalidade ).
Na região Nordeste, a rede privada concentrou 156 dos 201 cursos participantes, o equivalente a 77,6 \% do total da região.
Entre os cursos oferecidos por IES Privadas no Nordeste, 73 cursos foram avaliados com conceito 2, o conceito modal.
Dez cursos receberam os conceitos 1 e outros dez, conceito 4.
A 56 cursos foi atribuído conceito 3, e dois cursos ficaram com conceito 5.
Nessa combinação de categoria e região, cinco cursos ficaram Sem Conceito ( SC ).
As Instituições Públicas da região Nordeste participaram com 45 cursos ( 22,4 \% do total da região ).
Desses, 18 obtiveram conceito 3, o valor modal, seguidos por outros 13 cursos que receberam conceito 2.
Dois cursos receberam conceito 1 e sete cursos receberam conceito 4, e o conceito 5 foi atribuído a quatro cursos.
Um curso ficou Sem Conceito ( SC ) nessa região.
A maioria dos cursos na região Nordeste foi oferecida na Modalidade de Ensino Presencial : 194 cursos.
Para esses, a moda foi o conceito 2, com 83 cursos.
Os demais cursos presenciais da região foram alocados aos conceitos 1 ( 12 cursos ), 3 ( 72 cursos ), 4 ( 16 cursos ) e 5 ( seis cursos ).
Ficaram Sem Conceito ( SC ), cinco cursos.
Já na modalidade à Distância, três cursos receberam conceito 2, dois cursos, conceito 3 e um curso, conceito 4.
Um curso ficou Sem Conceito ( SC ).
Na região Sudeste, a proporção de cursos da rede privada, 93,3 \%, foi mais elevada do que a registrada nas demais regiões brasileiras, correspondendo a 431 dos 462 cursos participantes.
Nessa categoria e região, o conceito modal foi 3, atribuído a 202 cursos. 11 cursos ficaram Sem Conceito ( SC ) e os demais foram avaliados com os conceitos 1 ( 13 cursos ), 2 ( 114 cursos ), 4 ( 68 cursos ) e 5 ( 23 cursos ).
Dos 31 cursos oferecidos por Instituições Públicas na região Sudeste, o conceito modal foi 4, atribuído a 13 cursos.
Nessa região, nenhum curso de IES Públicas ficou Sem Conceito ( SC ).
Os demais cursos foram avaliados com os conceitos 1 ( um curso ), 2 ( cinco cursos ), 3 ( sete cursos ) e 5 ( cinco cursos ).
Também na região Sudeste, a maioria dos cursos foi oferecido na Modalidade de Ensino Presencial : 440 cursos.
Para esses, a moda foi o conceito 3, com 198 cursos.
Os demais cursos presenciais da região foram atribuídos aos conceitos 1 ( 14 cursos ), 2 ( 114 cursos ), 4 ( 76 cursos ) e 5 ( 27 cursos ).
Ficaram Sem Conceito ( SC ) 11 cursos.
Na região Sudeste, os 22 cursos oferecidos na Modalidade de Ensino Educação a Distância obtiveram o conceito 3 como conceito modal, atribuído a 11 cursos.
Os demais cursos foram avaliados com os conceitos 2 ( cinco cursos ), 4 ( cinco cursos ) e 5 ( um curso ).
Para essa Modalidade de Ensino e região, nenhum curso ficou Sem Conceito ( SC ) ou recebeu conceito 1.
As Instituições Privadas concentraram 206 dos 236 cursos participantes da região Sul, 87,3 \% do total regional.
Desses, 113 cursos obtiveram conceito 3, o conceito modal.
Nesta combinação de Categoria Administrativa e Grande Região, 54 cursos receberam conceito 2, 31 cursos receberam conceito 4, um curso recebeu conceito 5, e quatro cursos foram avaliados com conceito 1.
Nesta combinação de Categoria Administrativa e Grande Região, três cursos ficaram Sem Conceito ( SC ).
As Instituições Públicas na região Sul participaram com 30 cursos ( 12,7 \% dos cursos da região ), aos quais foram atribuídos os conceitos 2 ( dois cursos ), 3 ( 13 cursos, conceito modal ), 4 ( 11 cursos ) e 5 ( quatro cursos ).
Nenhum curso da região Sul, oferecido por IES Públicas, ficou Sem Conceito ( SC ) ou recebeu conceito 1.
Dos 236 cursos oferecidos na região Sul, a maioria dos cursos foi na Modalidade de Ensino Presencial : 223 cursos.
Para esses, a moda foi o conceito 3, com 122 cursos.
Os demais cursos presenciais da região foram atribuídos aos conceitos 1 ( quatro cursos ), 2 ( 49 cursos ), 4 ( 40 cursos ) e 5 ( cinco cursos ).
Três cursos ficaram Sem Conceito ( SC ).
Na região Sul, os 13 cursos oferecidos na Modalidade de Ensino Educação a Distância obtiveram o conceito 2, o conceito modal, atribuído a sete cursos.
Os demais cursos foram avaliados com os conceitos 3 ( quatro cursos ) e 4 ( dois cursos ).
Para essa Modalidade de Ensino e região, nenhum curso ficou Sem Conceito ( SC ) ou recebeu os conceitos 1 ou 5.
Na região Centro-Oeste, 102 dos 125 cursos participantes eram de Instituições Privadas ( 81,6 \% em termos regionais ).
Desses, 46 receberam conceito 2, o conceito modal.
Dos demais cursos, 31 receberam conceito 3, dez cursos, o conceito 4.
Um curso recebeu conceito 5, e 11 cursos foram avaliados com conceito 1.
Nesta combinação de Categoria Administrativa e Grande Região, três cursos ficaram Sem Conceito ( SC ).
Quanto aos 23 cursos oferecidos por Instituições Públicas na região Centro-Oeste ( 18,4 \% do total regional ), 11 cursos apresentaram conceito 3, o conceito modal, e nenhum curso ficou Sem Conceito ( SC ).
Os demais cursos foram avaliados com os conceitos 5 ( dois cursos ), 4 ( quatro cursos ), três cursos receberam conceito 1, e outros três receberam conceito 2.
Dos 125 cursos oferecidos na região Centro-Oeste, 120 cursos foram oferecidos na Modalidade de Ensino Presencial.
Para esses, a moda foi o conceito 2, com 47 cursos.
Os demais cursos Presenciais da região foram atribuídos aos conceitos 1 ( 14 cursos ), 3 ( 40 cursos ), 4 ( 13 cursos ) e 5 ( três cursos ).
Ficaram Sem Conceito ( SC ) três cursos.
Na região Centro-Oeste, dos cinco cursos oferecidos na Modalidade de Ensino Educação a Distância, nenhum curso ficou Sem Conceito ( SC ) ou recebeu os conceitos 1 ou 5.
Os demais cursos foram alocados aos conceitos 2 ( dois cursos ), 3 ( dois cursos ) e 4 ( um curso ).
No Gráfico 5.2, apresenta -se a distribuição do Conceito Enade, segundo a Categoria Administrativa da IES.
Os cursos em IES Públicas ( linha azul ) apresentam uma poligonal mais à direita do que os em IES Privadas ( linha verde ) e, consequentemente, uma distribuição de Conceitos Enade com valores maiores. 5.3 CONCEITOS POR ORGANIZAÇÃO ACADÊMICA E GRANDE REGIÃO Na Tabela 5.3, encontra -se a distribuição dos conceitos atribuídos aos cursos participantes do Enade/2018, na Área de Ciências Contábeis, por Organização Acadêmica, segundo as Grandes Regiões brasileiras.
Dos 1.101 cursos de Ciências Contábeis participantes, 350 eram oferecidos em Universidades, 220, em Centros Universitários, 529, em Faculdades e 2, em CEFET/IFET.
Esta distribuição corresponde a, respectivamente, 31,8 \%, 20,0 \%, 48,0 \% e 0,2 \% dos cursos.
De acordo com os dados apresentados, de todos os 42 cursos avaliados com conceito 5, 22 eram vinculados a Universidades.
Os cursos vinculados a esse tipo de Organização Acadêmica tiveram o conceito 3 como conceito modal, atribuído a 170 cursos ( 48,6 \% ).
Os demais cursos vinculados a Universidades receberam os conceitos 1 ( seis cursos ), 2 ( 74 cursos ), 4 ( 75 cursos ) e 5 ( 22 cursos, como já mencionado ).
Ficaram Sem Conceito ( SC ) três cursos.
Dos cursos em Centros Universitários, o conceito modal também foi 3, atribuído a 108 cursos ( 49,1 \% ).
Três cursos vinculados a esse tipo de Organização Acadêmica ficaram Sem Conceito ( SC ), e os demais cursos receberam os conceitos 1 ( dois cursos ), 2 ( 60 cursos ), 4 ( 42 cursos ) e 5 ( cinco cursos ).
Dos 529 cursos mantidos por Faculdades, 214 ( 40,5 \% ) receberam conceito 2, o conceito modal.
O segundo conceito mais frequentemente atribuído aos cursos oferecidos por essa categoria de IES foi o conceito 3, recebido por 198 cursos.
Na sequência, 43 cursos obtiveram conceito 4, 42 cursos, conceito 1 e 15 cursos, conceito 5.
Dos cursos oferecidos em Faculdades, 17 ficaram Sem Conceito ( SC ).
Os dois cursos mantidos por CEFET/IFET receberam conceito 3.
Considerando -se separadamente as regiões brasileiras, verifica -se que, na região Norte, as Universidades participaram com 18 dos 77 cursos avaliados, ao qual foi atribuído o conceito 3 como modal.
Os demais cursos foram atribuídos aos conceitos 1 ( um curso ), 2 ( quatro cursos ) e 4 ( três cursos ).
Nenhum curso ficou Sem Conceito ( SC ) ou recebeu conceito 5.
Os Centros Universitários da região Norte foram representados por 12 cursos, aos quais foram atribuídos os conceitos 2 ( sete cursos, conceito modal ), 3 ( três cursos ) e 4 ( dois cursos ).
Nenhum curso ficou Sem Conceito ( SC ) ou recebeu os conceitos 1 ou 5, nesta combinação de categoria e região.
As Faculdades participaram com 47 cursos na região Norte, que foram avaliados com os conceitos 1 ( cinco cursos ), 3 ( 14 cursos ) e 4 ( um curso ).
O conceito modal para essa categoria e região foi 2, com 27 cursos.
Nenhum curso ficou Sem Conceito ( SC ) ou recebeu conceito 5.
Já os CEFET/IFET não participaram com cursos na região Norte.
Na região Nordeste, as Universidades participaram com 59 dos 201 cursos da Área de Ciências Contábeis oferecidos.
Apenas um desses cursos ficou Sem Conceito ( SC ).
O conceito modal 3 foi atribuído a 27 cursos.
Os demais cursos, obtiveram os conceitos 1 ( dois cursos ), 2 ( 18 cursos ), 4 ( sete cursos ) e 5 ( quatro cursos ).
Os Centros Universitários contaram com 36 cursos participantes na região Nordeste.
Um curso ficou Sem Conceito ( SC ) e nenhum curso recebeu os conceitos 1 ou 5, nesse tipo de Organização Acadêmica.
O conceito 2 foi atribuído a 14 cursos, e outros quatro receberam o conceito 3 e o conceito 4 foi atribuído a sete cursos.
As Faculdades foram representadas por 106 cursos na região Nordeste, 54 dos quais receberam conceito 2, o conceito modal.
Os demais cursos foram avaliados com os conceitos 1 ( dez cursos ), 3 ( 33 cursos ), 4 ( três cursos ) e 5 ( dois cursos ).
Para esse tipo de Organização Acadêmica e região, quatro cursos ficaram Sem Conceito ( SC ).
Os CEFET/IFET não participaram com cursos na região Nordeste.
Na região Sudeste, as Universidades concentraram 138 dos 462 cursos de Ciências Contábeis da região.
Dos cursos oferecidos em Universidades nessa região, o conceito modal foi 3, atribuído a 66 cursos.
Os demais cursos receberam os conceitos 1 ( um curso ), 2 ( 23 cursos ), 4 ( 36 cursos ) e 5 ( 11 cursos ).
Um curso ficou Sem Conceito ( SC ).
Os Centros Universitários participaram com 113 cursos na região Sudeste, 60 dos quais obtiveram o conceito modal 3.
O conceito 1 foi atribuído a um curso, o conceito 2, a 28 cursos.
O conceito 4 foi atribuído a 18 cursos e o conceito 5, a cinco cursos.
Um curso ficou Sem Conceito ( SC ) nessa combinação de Organização Acadêmica e região.
Já as Faculdades foram representadas por 211 cursos na região Sudeste, das quais nove ficaram Sem Conceito ( SC ).
O conceito modal 3 foi atribuído a 83 cursos.
Aos demais cursos foram atribuídos os conceitos 1 ( 12 cursos ), 2 ( 68 cursos ), 4 ( 27 cursos ) e 5 ( 12 cursos ).
Os CEFET/IFET não participaram com cursos na região Sudeste.
Dos 236 cursos da região Sul, 104 eram mantidos por Universidades, para os quais o conceito modal foi 3, atribuído a 54 cursos.
Os demais cursos receberam os conceitos 1 ( um curso ), 2 ( 21 cursos ), 4 ( 22 cursos ) e 5 ( cinco cursos ).
Nesse tipo de organização, um curso da região Sul ficou Sem Conceito ( SC ).
Dos 38 cursos participantes de Centros Universitários da região Sul, nenhum curso ficou Sem Conceito ( SC ) ou recebeu os conceitos 1 ou 5.
Os conceitos 2, 3 e 4 foram atribuídos, respectivamente, a seis cursos, a 20 cursos e a 12 cursos.
Quanto aos 92 cursos vinculados a Faculdades na região Sul, 50 receberam o conceito modal 3.
O conceito 2 foi atribuído a 29 cursos, o conceito 4, a oito cursos, já o conceito 1, a três cursos.
Dois cursos ficaram Sem Conceito ( SC ) ou recebeu conceito 5.
Já os CEFET/IFET participaram com dois cursos na região Sul, avaliados com conceito 3.
Na região Centro-Oeste, 31 dos 125 cursos oferecidos eram mantidos por Universidades.
Para este tipo de Organização Acadêmica, nenhum curso ficou Sem Conceito ( SC ), e a 13 cursos foi atribuído conceito modal, 3.
Os demais cursos foram atribuídos aos conceitos 1 ( um curso ), 2 ( oito cursos ), 4 ( sete cursos ) e 5 ( dois cursos ).
Os Centros Universitários da região Centro-Oeste contaram com 21 cursos, 11 dos quais foram avaliados com conceito 3, o conceito modal.
Os demais cursos receberam os conceitos 1 ( um curso ), 2 ( cinco cursos ), e 4 ( três cursos ).
Um curso ficou Sem Conceito ( SC ) ou recebeu conceito 5.
Dos 73 cursos oferecidos em Faculdades na região Centro-Oeste, 36 receberam o conceito modal 2, e 12 obtiveram conceito 1.
Nesta região, dois cursos ficaram Sem Conceito ( SC ), e os conceitos 3, 4 e 5 foram atribuídos, respectivamente, a 18 cursos, a quatro cursos e a um curso.
Já os CEFET/IFET não participaram com cursos na região Centro-Oeste.
No Gráfico 5.3, apresenta -se a distribuição do Conceito Enade, segundo a Organização Acadêmica da IES.
Os cursos em CEFET/IFET ( linha laranja ) \textbf{estão} concentrados no conceito 2.
Não se levando em conta os cursos em CEFET/IFET, os cursos em Faculdades ( linha vermelha ), são representados por uma poligonal mais à esquerda e, consequentemente, distribuição de Conceitos Enade com menores valores.
No outro extremo, com uma distribuição semelhante entre si, encontram -se os cursos mantidos por Universidades ( linha azul ) e Centros Universitários ( linha verde ), apresentando as respectivas poligonais mais à direita e, consequentemente, distribuição de Conceitos Enade com maiores valores.

% matched lemmas: cada, estar
\end{textsample}
